%%% The main file. It contains definitions of basic parameters and includes all other parts.

%% Settings for single-side (simplex) printing
% Margins: left 40mm, right 25mm, top and bottom 25mm
% (but beware, LaTeX adds 1in implicitly)
\documentclass[12pt,a4paper]{report}
\setlength\textwidth{145mm}
\setlength\textheight{247mm}
\setlength\oddsidemargin{15mm}
\setlength\evensidemargin{15mm}
\setlength\topmargin{0mm}
\setlength\headsep{0mm}
\setlength\headheight{0mm}
% \openright makes the following text appear on a right-hand page
\let\openright=\clearpage

%% Settings for two-sided (duplex) printing
% \documentclass[12pt,a4paper,twoside,openright]{report}
% \setlength\textwidth{145mm}
% \setlength\textheight{247mm}
% \setlength\oddsidemargin{14.2mm}
% \setlength\evensidemargin{0mm}
% \setlength\topmargin{0mm}
% \setlength\headsep{0mm}
% \setlength\headheight{0mm}
% \let\openright=\cleardoublepage

%% Generate PDF/A-2u
\usepackage[a-2u]{pdfx}

%% Character encoding: usually latin2, cp1250 or utf8:
\usepackage[utf8]{inputenc}

%% Prefer Latin Modern fonts
\usepackage{lmodern}

%% Further useful packages (included in most LaTeX distributions)
\usepackage{amsmath}        % extensions for typesetting of math
\usepackage{amsfonts}       % math fonts
\usepackage{amsthm}         % theorems, definitions, etc.
\usepackage{bbding}         % various symbols (squares, asterisks, scissors, ...)
\usepackage{bm}             % boldface symbols (\bm)
\usepackage{graphicx}       % embedding of pictures
\usepackage{fancyvrb}       % improved verbatim environment
\usepackage{natbib}         % citation style AUTHOR (YEAR), or AUTHOR [NUMBER]
\usepackage[nottoc]{tocbibind} % makes sure that bibliography and the lists
			    % of figures/tables are included in the table
			    % of contents
\usepackage{dcolumn}        % improved alignment of table columns
\usepackage{booktabs}       % improved horizontal lines in tables
\usepackage{paralist}       % improved enumerate and itemize
\usepackage{xcolor}         % typesetting in color
\usepackage{tikz}
\usetikzlibrary{positioning}
\usetikzlibrary{arrows}
%%% Basic information on the thesis

% Thesis title in English (exactly as in the formal assignment)
\def\ThesisTitle{Fitness and novelty in evolutionary reinforcement learning}

% Author of the thesis
\def\ThesisAuthor{František Dostál}

% Year when the thesis is submitted
\def\YearSubmitted{2023}

% Name of the department or institute, where the work was officially assigned
% (according to the Organizational Structure of MFF UK in English,
% or a full name of a department outside MFF)
\def\Department{Name of the department}

% Is it a department (katedra), or an institute (ústav)?
\def\DeptType{Katedra teoretické informatiky a matematické logiky}

% Thesis supervisor: name, surname and titles
\def\Supervisor{Mgr. Roman Neruda, CSc.}

% Supervisor's department (again according to Organizational structure of MFF)
\def\SupervisorsDepartment{department}

% Study programme and specialization
\def\StudyProgramme{Mgr.}
\def\StudyBranch{Artificial Inteligence}

% An optional dedication: you can thank whomever you wish (your supervisor,
% consultant, a person who lent the software, etc.)
\def\Dedication{%
Dedication.
}

% Abstract (recommended length around 80-200 words; this is not a copy of your thesis assignment!)
\def\Abstract{%
Novelty is a novel aproach to modeling selection criteria in evolutionary algorithms and has been proven as viable technique of avoiding pitfalls of false optima in tasks abundant with them, such as solving mazes. Rather than closing the topic however, this finding opened other problems to explore: How to properly apply novelty in tasks that yield slightely better to conventional aproaches? How to properly model behavioral space nessessary for novelty computation? In this thesis we investigate use of novelty in selected reinforcement learnig tasks, it's combinations with classical fitness and propose behavior space models for the respective RL tasks.
}

% 3 to 5 keywords (recommended), each enclosed in curly braces
\def\Keywords{%
{evolution} {novelty} {fitness} {behavioral space}
}

%% The hyperref package for clickable links in PDF and also for storing
%% metadata to PDF (including the table of contents).
%% Most settings are pre-set by the pdfx package.
\hypersetup{unicode}
\hypersetup{breaklinks=true}

% Definitions of macros (see description inside)
%%% This file contains definitions of various useful macros and environments %%%
%%% Please add more macros here instead of cluttering other files with them. %%%

%%% Minor tweaks of style

% These macros employ a little dirty trick to convince LaTeX to typeset
% chapter headings sanely, without lots of empty space above them.
% Feel free to ignore.
\makeatletter
\def\@makechapterhead#1{
  {\parindent \z@ \raggedright \normalfont
   \Huge\bfseries \thechapter. #1
   \par\nobreak
   \vskip 20\p@
}}
\def\@makeschapterhead#1{
  {\parindent \z@ \raggedright \normalfont
   \Huge\bfseries #1
   \par\nobreak
   \vskip 20\p@
}}
\makeatother

% This macro defines a chapter, which is not numbered, but is included
% in the table of contents.
\def\chapwithtoc#1{
\chapter*{#1}
\addcontentsline{toc}{chapter}{#1}
}

% Draw black "slugs" whenever a line overflows, so that we can spot it easily.
\overfullrule=1mm

%%% Macros for definitions, theorems, claims, examples, ... (requires amsthm package)

\theoremstyle{plain}
\newtheorem{thm}{Theorem}
\newtheorem{lemma}[thm]{Lemma}
\newtheorem{claim}[thm]{Claim}

\theoremstyle{plain}
\newtheorem{defn}{Definition}

\theoremstyle{remark}
\newtheorem*{cor}{Corollary}
\newtheorem*{rem}{Remark}
\newtheorem*{example}{Example}

%%% An environment for proofs

\newenvironment{myproof}{
  \par\medskip\noindent
  \textit{Proof}.
}{
\newline
\rightline{$\qedsymbol$}
}

%%% An environment for typesetting of program code and input/output
%%% of programs. (Requires the fancyvrb package -- fancy verbatim.)

\DefineVerbatimEnvironment{code}{Verbatim}{fontsize=\small, frame=single}

%%% The field of all real and natural numbers
\newcommand{\R}{\mathbb{R}}
\newcommand{\N}{\mathbb{N}}

%%% Useful operators for statistics and probability
\DeclareMathOperator{\pr}{\textsf{P}}
\DeclareMathOperator{\E}{\textsf{E}\,}
\DeclareMathOperator{\var}{\textrm{var}}
\DeclareMathOperator{\sd}{\textrm{sd}}

%%% Transposition of a vector/matrix
\newcommand{\T}[1]{#1^\top}

%%% Various math goodies
\newcommand{\goto}{\rightarrow}
\newcommand{\gotop}{\stackrel{P}{\longrightarrow}}
\newcommand{\maon}[1]{o(n^{#1})}
\newcommand{\abs}[1]{\left|{#1}\right|}
\newcommand{\dint}{\int_0^\tau\!\!\int_0^\tau}
\newcommand{\isqr}[1]{\frac{1}{\sqrt{#1}}}

%%% Various table goodies
\newcommand{\pulrad}[1]{\raisebox{1.5ex}[0pt]{#1}}
\newcommand{\mc}[1]{\multicolumn{1}{c}{#1}}


% Title page and various mandatory informational pages
\begin{document}

%%% Title page of the thesis and other mandatory pages

%%% Title page of the thesis

\pagestyle{empty}
\hypersetup{pageanchor=false}
\begin{center}

\centerline{\mbox{\includegraphics[width=166mm]{../img/logo-en.pdf}}}

\vspace{-8mm}
\vfill

{\bf\Large MASTER THESIS}

\vfill

{\LARGE\ThesisAuthor}

\vspace{15mm}

{\LARGE\bfseries\ThesisTitle}

\vfill

\Department

\vfill

{
\centerline{\vbox{\halign{\hbox to 0.45\hsize{\hfil #}&\hskip 0.5em\parbox[t]{0.45\hsize}{\raggedright #}\cr
Supervisor of the master thesis:&\Supervisor \cr
\noalign{\vspace{2mm}}
Study programme:&\StudyProgramme \cr
\noalign{\vspace{2mm}}
Study branch:&\StudyBranch \cr
}}}}

\vfill

% Zde doplňte rok
Prague \YearSubmitted

\end{center}

\newpage

%%% Here should be a bound sheet included -- a signed copy of the "master
%%% thesis assignment". This assignment is NOT a part of the electronic
%%% version of the thesis. DO NOT SCAN.

%%% A page with a solemn declaration to the master thesis

\openright
\hypersetup{pageanchor=true}
\pagestyle{plain}
\pagenumbering{roman}
\vglue 0pt plus 1fill

\noindent
I declare that I carried out this master thesis independently, and only with the cited
sources, literature and other professional sources. It has not been used to obtain another
or the same degree.

\medskip\noindent
I understand that my work relates to the rights and obligations under the Act No.~121/2000 Sb.,
the Copyright Act, as amended, in particular the fact that the Charles
University has the right to conclude a license agreement on the use of this
work as a school work pursuant to Section 60 subsection 1 of the Copyright~Act.

\vspace{10mm}

\hbox{\hbox to 0.5\hsize{%
In \hbox to 6em{\dotfill} date \hbox to 6em{\dotfill}
\hss}\hbox to 0.5\hsize{\dotfill\quad}}
\smallskip
\hbox{\hbox to 0.5\hsize{}\hbox to 0.5\hsize{\hfil Author's signature\hfil}}

\vspace{20mm}
\newpage

%%% Dedication

\openright

\noindent
\Dedication

\newpage

%%% Mandatory information page of the thesis

\openright

\vbox to 0.5\vsize{
\setlength\parindent{0mm}
\setlength\parskip{5mm}

Title:
\ThesisTitle

Author:
\ThesisAuthor

\DeptType:
\Department

Supervisor:
\Supervisor, \SupervisorsDepartment

Abstract:
\Abstract

Keywords:
\Keywords

\vss}

\newpage

\openright
\pagestyle{plain}
\pagenumbering{arabic}
\setcounter{page}{1}


%%% A page with automatically generated table of contents of the master thesis

\tableofcontents

%%% Each chapter is kept in a separate file
\chapter*{Introduction}
\addcontentsline{toc}{chapter}{Introduction}


\chapter{Introduction to Artificial Agents}
To understand the reasoning guiding this work we have to first look at the fundamental subject of our efforts - the artificial agents.
\begin{defn}\citep{russell2021}
A rational agent is one that acts so as to achieve
best outcome or, when there is uncertainty, the best expected outcome.
\end{defn}
\section{Reactive agents}



\section{Reinforcement learning}
\begin{defn}
A reward is a numerical evaluation of agent's behavior within a given enviroment. A reward function then is a function which assigns pair of state and action a real numerical value.
\end{defn}
\chapter{Evolutionary algorithms}

\section{•}

\section{Numerical Algorithms}
In this section we will be reviewing different aproaches to implementation of evolutionary algoAlgorithmsrithms which were considered
for the purpose of this thesis.
Because we are optimising
The algorithms were selected based on usefulness for the purposes of numerical optimalisation but also for their representativness of certain aspects of the field.
Since the aim of this work is to compare aproaches of formulating utility function optimised by the evolution rather than the algorithms as such, 
we will consider them in their simplest form still expedient for satisfactory optimalisation of selected tasks.
\subsection{Evolutionary Strategies}
Evolutionary strategies are 
\subsection{Differential Evolution}


\section{Utility functions}
\subsection{Fitness}
\subsection{Pure Novelty}
\subsection{Combinations of Novelty}

\chapter{Enviroments}
\section{Gymnasium}
\subsection{CartPole}
\subsection{Moon Landing}
\chapter{Evolutionary algorithms}
While reinforcement learning is a useful method of developing expert agents, it has certain drawbacks. Since it only iterates on handful of agents (usually a single agent) with the feedback directly recieved, it does not search the agent space very thoroughly and the resulting adapted agent can be ill-prepared for conditions arising from its good performance in the previous stages of more structured enviroments. 
If this might be a concern, more suitable way of producing expert agents, might be an evolutionary algorithm. There the agent is typically evaluated on the complex of its behavior within the enviroment.
\begin{defn}
An objective function is a mathematical expression that defines the goal of an optimization problem, representing the quantity that needs to be maximized or minimized.
\end{defn}
The evolutionary algorithms are a class of optimalisation methods modeled after our idea of biological process of evolution. Subjects of these optimalisation algorithms are called individuals, often representing feasible solutions of the optimalisation problem. The evolutionary algorithms are specific in that they search the space of individuals semi-randomly using folowing operations:
\begin{itemize}
	\item recombination (or cross-over)
	\item mutation
	\item selection 
\end{itemize}
The selection process usually utilises some kind of value assigned to an individual. Because the purpose of this thesis is to compare performance of different ways to assign such value, some newer than others, we have to sidestep the established terminology in vast majority of the literature a little bit.
\begin{defn}
Let $ A $ be an evolutionary algorithm with objective function $ g: I\rightarrow R $ searching over space of individuals $ I $. Then any function $ f: I\rightarrow R $ used to guide selection in A would be called utility function. Only in the special case where $ g=f $ we will call it fitness function.
\end{defn}
\section{Utility functions}

\subsection{Fitness}
Using objective function as the utility function of the algorithm is the most direct way of establishing relationship between the solution we want and the individuals in the evolving population.\\
Though this aproach ensures that the most optimal individuals from given population  will generally be selected into the next generation during the evolution it does not guarantee that their potential for evolving is not diminished either. This way the evolution using fitness can easily get stuck in local optimum of the objective function. While in most cases it might be still possible to reach global optimum on the given problem by changing the hyperparameters to favor random exploration, it may come to point where the benefits of evolutionary aproach are made void.
\subsection{Pure Novelty}
In the paper \cite{noveltyarticle} a new aproach to utility functions has been proposed, called novelty, designed to help avoid traps of local optima. 
\begin{defn}\citep{doncieux:hal-02561846}
Let E be enviroment and S be a space of all possible states of the enviroment and let $ d(x,y) $ be some metric. Then (B, d) is a behavioral space and a function $ o_B: S^{T}\rightarrow B $ is called observer function. 
\end{defn}
\begin{defn}\citep{noveltyarticle}
Let P be a population of individuals within EA and (Img(B), d) be a behavioral space. Then for each individual $ a $ we define novelty as such:  $$ N(a) =\frac{\sum_{b\in P, b\neq a}d(a,b)}{|P|-1} $$ 
\end{defn}
With novelty as utility function the selection process incentivises behaviours inversely to their representation in the current population, in other words according to how "novel" they are. This allows the evolutionary search to retain variety within the population across generations, leveraging the recombination as well as mutations to leave local optima.\\
Downside of this aproach is in the ambiguity in the choice of behavioral space. Incorporating too many enviroment variables will lead to the curse of dimensionality. Leaving out an important variable means the algorithm will not retain diversity where it matters, making novelty pointless. This means the behavioral function and space have to be selected empirically for given problem.
\subsection{Combinations of Novelty}
The novelty does not take into account the individual's performce with objective function. Therefore when two invidials with similair behavioral characteristics appear but one of them performs significantly better than the other, pure novelty algorithm can easily select the less performant individual, losing the more performant genom to further generations. To prevent this we want to meaningfully combine novlety with fitness.\\

\section{Continuous Optimalisation}
\subsection{Evolutionary Strategies \citep{evoStrategy}}
Evolutionary strategies are 
\subsection{Differential Evolution}


\chapter{Comparison criteria}
\chapter{Experiments}
\chapter*{Conclusion}
\addcontentsline{toc}{chapter}{Conclusion}


%%% Bibliography
%%% Bibliography (literature used as a source)
%%%
%%% We employ bibTeX to construct the bibliography. It processes
%%% citations in the text (e.g., the \cite{...} macro) and looks up
%%% relevant entries in the bibliography.bib file.
%%%
%%% The \bibliographystyle command selects, which style will be used
%%% for references from the text. The argument in curly brackets is
%%% the name of the corresponding style file (*.bst). Both styles
%%% mentioned in this template are included in LaTeX distributions.

\bibliographystyle{plainnat}    %% Author (year)
% \bibliographystyle{unsrt}     %% [number]

\renewcommand{\bibname}{Bibliography}

%%% Generate the bibliography. Beware that if you cited no works,
%%% the empty list will be omitted completely.

\bibliography{bibliography}

%%% If case you prefer to write the bibliography manually (without bibTeX),
%%% you can use the following. Please follow the ISO 690 standard and
%%% citation conventions of your field of research.

% \begin{thebibliography}{99}
%
% \bibitem{lamport94}
%   {\sc Lamport,} Leslie.
%   \emph{\LaTeX: A Document Preparation System}.
%   2nd edition.
%   Massachusetts: Addison Wesley, 1994.
%   ISBN 0-201-52983-1.
%
% \end{thebibliography}


%%% Figures used in the thesis (consider if this is needed)
\listoffigures

%%% Tables used in the thesis (consider if this is needed)
%%% In mathematical theses, it could be better to move the list of tables to the beginning of the thesis.
\listoftables

%%% Abbreviations used in the thesis, if any, including their explanation
%%% In mathematical theses, it could be better to move the list of abbreviations to the beginning of the thesis.
\chapwithtoc{List of Abbreviations}

%%% Attachments to the master thesis, if any. Each attachment must be
%%% referred to at least once from the text of the thesis. Attachments
%%% are numbered.
%%%
%%% The printed version should preferably contain attachments, which can be
%%% read (additional tables and charts, supplementary text, examples of
%%% program output, etc.). The electronic version is more suited for attachments
%%% which will likely be used in an electronic form rather than read (program
%%% source code, data files, interactive charts, etc.). Electronic attachments
%%% should be uploaded to SIS and optionally also included in the thesis on a~CD/DVD.
%%% Allowed file formats are specified in provision of the rector no. 72/2017.
\appendix
\chapter{Attachments}

\section{First Attachment}

\openright
\end{document}
